\section{Modelo Intermunicipal Sem Intermediação Estadual}

Este modelo representa um cenário hipotético de operação descentralizada, na qual no qual municípios podem redistribuir vacinas diretamente entre si sem a necessidade de passagem por órgãos intermediários, como secretarias estaduais de saúde, conforme representado a seguir:

\[
\text{Município } i \;\longrightarrow\; \text{Município } j
\]

A adoção dessa abordagem tem como principal objetivo simplificar a estrutura do modelo e, ao mesmo tempo, possibilitar avaliar o desempenho, em termos de tempo, de uma estratégia de distribuição direta.

\subsection{Variável de Decisão e Parâmetros}

A \textbf{variável de decisão} é definida:
\\
\\
\begin{equation*}
    x_{ijv} \in \mathbb{Z}_{\geq 0} \quad \forall i \in I, j \in J, t \in V,
\end{equation*}
onde:
\begin{itemize}
    \item $I$ representa os municípios $i$ mapeados;
    \item $J$ representa os municípios $j$ mapeados;
    \item $v$ representa os tipos de vacina mapeados.
    \item $x$ representa a quantidade de vacinas transportadas;
    \item $i$ indica o município de origem $i$, que possui oferta da vacina que será transportada;
    \item $j$ indica o município de destino $j$, que possui demanda pela vacina que será transportada;
    \item $v$ indica o tipo de vacina que será transportada, ou seja, definido por uma tabela indicando o índice apropriado para cada tipo. Por exemplo:
\end{itemize}

\begin{quadro}[htb] 
\centering
\caption{Exemplo de uso da variável $v$}
\label{quad:quadro1}
% Trocamos tabularx por tabular e removemos a largura fixa
\begin{tabular}{|c|l|} 
\hline
\textbf{$v$} & \textbf{Nome da vacina} \\
\hline
1 & BCG \\
\hline
2 & HEPATITE B \\
\hline
\end{tabular}
\end{quadro}

Além disso, como x representa a quantidade transportada, foi definido que seu valor deve ser um número inteiro e não-negativo, pois não faria sentido transportar meia dose de vacina, por exemplo. Desta forma, se trata de um \textbf{Problema de Programação Linear Inteira}.

Além da variável de decisão, é importante definir alguns \textbf{parâmetros}.

As principais opções reais de modais para o transporte de vacinas são o rodoviário e o aéreo e por isso eles serão considerados na modelagem. Esses modais possuem tempos de deslocamento diferentes para uma mesma distância, por isso se faz necessário definir esses \textbf{parâmetros de tempo}.

Sejam:

\[
\tau_{ij}^{rod}, \quad \tau_{ij}^{areo} \in \mathbb{R}_{\geq 0}
\]

onde:
\begin{itemize}
    \item $\tau_{ij}^{rod}$ é o tempo de deslocamento do modal rodoviário entre o município de origem $i$ e o município de destino $j$.
    \item $\tau_{ij}^{areo}$ é o tempo de deslocamento do modal aéreo entre o município de origem $i$ e o município de destino $j$.
\end{itemize}

No caso do transporte aéreo, o modelo não considera explicitamente apenas o tempo de voo entre os municípios. O tempo $\tau_{ij}^{areo}$ representa o tempo total agregado da operação logística aérea, incluindo:

\begin{itemize}
    \item transporte rodoviário do município de origem até o aeroporto;
    \item tempo de voo entre aeroportos;
    \item transporte rodoviário do aeroporto de destino até o município final.
\end{itemize}

Dessa forma, $\tau_{ij}^{areo}$ já incorpora implicitamente todas as etapas necessárias para a entrega via modal aéreo, mantendo o modelo compacto sem perda significativa de realismo.

Na prática, o processo de distribuição de vacinas não se limita apenas ao tempo de deslocamento entre os pontos de origem e destino. Ao longo da cadeia logística, existem etapas intermediárias que incluem atividades como desembarque, descarga, conferência, armazenamento temporário e redistribuição das vacinas.

Essas atividades introduzem um tempo adicional no processo, denominado neste trabalho como \textbf{tempo de processamento logístico} e tal tempo ocorre, principalmente, nos pontos de destino, como aeroportos, secretarias municipais de saúde ou centros de distribuição locais.

O tempo de processamento logístico pode variar de acordo com o modal de transporte utilizado. No caso do transporte aéreo, por exemplo, há etapas adicionais como operações aeroportuárias, inspeções e manuseio especializado, que tendem a aumentar o tempo total de processamento. Já no transporte rodoviário, o fluxo tende a ser mais direto, resultando em tempos menores. Para capturar esse comportamento, são considerados tempos de processamento logístico distintos por modal.

Modelar explicitamente essas etapas exigiria a introdução de novas variáveis de decisão e restrições, caracterizando um problema mais complexo e poderia comprometer a tratabilidade computacional. Em vez disso, os tempos de deslocamento e de processamento logístico e a escolha do modal mais adequado (rodoviário ou aéreo) serão incorporados, implicitamente, a \textbf{parâmetros de custo total}.

Sejam:

\[
\delta^{rod}, \quad \delta^{areo} \in \mathbb{R}_{\geq 0}
\]

onde:
\begin{itemize}
    \item $\delta^{rod}$ representa o tempo de processamento logístico do modal rodoviário.
    \item $\delta^{areo}$ representa o tempo de processamento logístico do modal aéreo.
\end{itemize}

Ainda, é necessário considerar as diferenças de custo financeiro e de complexidade logística entre os dois tipos de modais. O transporte aéreo é mais rápido, porém mais caro financeiramente e mais complexo logisticamente.

Seja:

\[
\alpha > 1
\]

um \textbf{fator de penalização associado ao transporte aéreo} e que reflete limitações operacionais e custos adicionais do transporte aéreo. 

Essa abordagem permite modelar a preferência pelo transporte rodoviário em situações onde ambos os modais são viáveis, ao mesmo tempo em que mantém o transporte aéreo como alternativa para cenários em que o ganho de tempo é significativo. Essa abordagem permite capturar as diferenças operacionais entre modais sem aumentar a complexidade do modelo.

Sejam:
\begin{itemize}
    \item $C_{ij}^{rod}$, o \textbf{Custo Total do transporte rodoviário};
    \item $C_{ij}^{areo}$, o \textbf{Custo Total do transporte aéreo}.
\end{itemize}


Dessa forma, o \textbf{Custo Total de transporte por modal} é dado por:

\[
C_{ij}^{rod} = \tau_{ij}^{rod} + \delta^{rod}
\]
\[
C_{ij}^{areo} = \alpha \cdot \tau_{ij}^{areo} + \delta^{areo}
\]


Dessa forma, define-se o \textbf{Custo Total de transporte} como:

\[
C_{ij} = \min \left( \quad C_{ij}^{rod} \quad , \quad C_{ij}^{areo} \quad \right),  \quad C_{ij} \in \mathbb{R}_{\geq 0}
\]

Assim, o custo total de transporte $C_{ij}$ é definido como o menor tempo efetivo entre os modais disponíveis.

Essa abordagem não introduz novas variáveis de decisão e, dessa forma, mantém a eficiência computacional do modelo, ao mesmo tempo em que aumenta sua aderência à realidade operacional do sistema de distribuição de vacinas.

Além das diferenças de tempo entre os modais de transporte, é necessário considerar que a \textbf{capacidade de transporte} também pode variar de acordo com o modal utilizado.

Sejam:
\[
u_{ij}^{\text{rod}}, \quad u_{ij}^{\text{areo}} \in \mathbb{R}_{\geq 0}
\]

onde:
\begin{itemize}
    \item $u_{ij}^{\text{rod}}$ representa a capacidade máxima de transporte entre os municípios $i$ e $j$ utilizando o modal rodoviário;
    \item $u_{ij}^{\text{areo}}$ representa a capacidade máxima de transporte entre os municípios $i$ e $j$ utilizando o modal aéreo.
\end{itemize}

Dessa forma, a capacidade efetiva de transporte entre os municípios $i$ e $j$ é definida de forma consistente com o custo escolhido, conforme:

\begin{equation}
u_{ij} =
\begin{cases}
u_{ij}^{\text{rod}}, & \text{se } C_{ij}^{\text{rod}} \le C_{ij}^{\text{areo}} \\
u_{ij}^{\text{areo}}, & \text{caso contrário}
\end{cases}
\end{equation}

Além disso, essa modelagem mantém a simplicidade do problema, evitando o aumento do número de variáveis de decisão, ao mesmo tempo em que incorpora diferenças operacionais relevantes entre os modais de transporte.

Com a variável de decisão e parâmetros definidos, pode-se definir o tempo de execução do plano de distribuição de vacinas, ou seja, quanto tempo levou até a última entrega chegar desde que a distribuição foi iniciada.

Considerando que as entregas acontecem em paralelo, então, o \textbf{tempo de distribuição} é dado por:


\[
T = \max \left\{ C_{ij} \;\middle|\; x_{ijv} > 0 \right\}
\]

onde:
\begin{itemize}
    \item $T$ é o tempo de distribuição das vacinas.
\end{itemize}

\subsection{Função Objetivo}

Assim, a \textbf{função objetivo} ficou definida da seguinte forma:

\begin{equation}
    \min z = \sum_{v} \sum_{i} \sum_{j} C_{ij} \cdot x_{ijv}
    \label{eq:intermunicipal_sem_penalidade}
\end{equation}

Portanto, o problema modelado tem como objetivo \textbf{minimizar o tempo necessário para entregar a quantidade necessária das doses de vacinas ofertadas para os municípios com demanda}.

Pode-se, ainda, fazer uma modificação da função anterior adicionando uma \textbf{penalidade $\rho_{iv}$} para a vacina de tipo $v$ do município $i$ de acordo com o tempo de expiração, para que, no processo de transporte, haja o menor número de doses expiradas possível. Assim, temos:


\[ 0 < \rho_{iv} \le 1, \]  \\
onde:
\begin{itemize}
    \item dado um tempo de expiração $e_{iv}$, ou seja, tempo para a vacina $v$ retida no município $i$ expirar, então a penalidade é definida por: $\rho_{iv} = \frac{1}{e_{iv}}$.
\end{itemize}

Assim, a função objetivo modificada fica definida da seguinte forma:

\begin{equation}
    \min z = \sum_{v} \sum_{i} \sum_{j} C_{ij} \cdot x_{ijv} \cdot \rho_{iv}
    \label{eq:intermunicipal_com_penalidade}
\end{equation}

Esta função tem um objetivo ligeiramente diferente: \textbf{minimizar o tempo ponderado pelo risco de expiração}.

Finalmente, algumas restrições precisam ser levadas em consideração para este problema.

\subsection{Restrições}

\textbf{A primeira restrição} é a \textbf{restrição de oferta} de cada vacina $v$ em cada município $i$. 
O número de vacinas transportadas pode ser menor que número de vacinas ofertadas no local de origem e, no máximo, pode-se transportar uma quantidade igual a quantidade ofertada pela origem. Assim, temos:

\begin{equation*}
    \sum_{j} x_{ijv} \le o_{iv}, \quad \forall i \in I, \; v \in V
\end{equation*}
onde:
\begin{itemize}
    \item $o_{ik}$ representa a quantidade de vacinas do tipo $v$ ofertadas no município $i$;
\end{itemize}

Seria um erro se a restrição fosse definida desta forma:
\begin{equation*}
    \sum_{j} x_{ijv} = o_{iv}, \quad \forall i \in I, \; v \in V
\end{equation*}
pois, nesse caso, todas as vacinas ofertadas seriam obrigadas a serem transportadas e, num cenário real, do conjunto de vacinas ofertadas pode sobrar vacina pode ou não haver destino viável.

\textbf{A segunda restrição} é a \textbf{restrição de demanda} da vacina $v$ por cada município $j$. 
O número de vacinas transportadas deve ser exatamente a quantidade demandada para que não haja falta ou desperdício. Assim, temos:

\begin{equation*}
    \sum_{i} x_{ijv} = d_{jv}, \quad \forall j \in J, \; v \in V
\end{equation*}
onde:
\begin{itemize}
    \item $d_{jv}$ representa a quantidade de vacinas do tipo $v$ demandadas no município $j$;
    \item $J$ representa os municípios $j$ mapeados.
\end{itemize}

E, nesse caso, estamos forçando o atendimento total da demanda.

Seria um erro se a restrição fosse definida desta forma:
\begin{equation*}
    \sum_{i} x_{ijv} \le d_{jv}, \quad \forall j \in J, \; v \in V
\end{equation*}
pois isso implicaria que se pode não atender a demanda e o modelo pode simplesmente não entregar nada (minimiza custo = 0).

\textbf{A terceira restrição} é a \textbf{restrição de capacidade} de entrega pelo veículo de transporte da origem $i$ ao destino $j$. 
O número de vacinas transportadas não pode exceder a capacidade do veículo que realiza o transporte. Assim, temos:

\begin{equation*}
    \sum_{v} x_{ijv} \le u_{ij}, \quad \forall i \in I, j \in J
\end{equation*}
onde:
\begin{itemize}
    \item $u_{ij}$ representa a capacidade máxima (limite superior) de vacinas que pode ser transportada do município $i$ para o município $j$.
\end{itemize}

Assume-se que diferentes tipos de vacina compartilham a mesma capacidade de transporte, o que é consistente com a operação logística real, na qual múltiplos imunizantes são transportados conjuntamente no mesmo veículo.

\textbf{A quarta restrição} é a \textbf{restrição de viabilidade temporal} e é relativa à ideia de remover pares onde o tempo de transporte é maior que o tempo restante até o vencimento. Assim, temos:

\begin{equation*}
x_{ijv} = 0, \quad \forall i \in I, \; j \in J, \; v \in V \text{ tais que } C_{ij} > e_{iv}
\end{equation*}

Em suma, o PPLI ficou definido assim:

\begin{equation*}
    \min z = \sum_{v} \sum_{i} \sum_{j} C_{ij} \cdot x_{ijv}
\end{equation*}
sujeito a:
\begin{equation*}
    \sum_{j} x_{ijv} \le o_{ik}, \quad \forall i \in I, \; v \in V
\end{equation*}

\begin{equation*}
    \sum_{i} x_{ijv} = d_{jv}, \quad \forall j \in J, \; v \in V
\end{equation*}

\begin{equation*}
    \sum_{v} x_{ijv} \le u_{ij}, \quad \forall i \in I, j \in J
\end{equation*}

\begin{equation*}
x_{ijv} = 0, \quad \forall i \in I, \; j \in J, \; v \in V \text{ tais que } C_{ij} > e_{iv}
\end{equation*}

\begin{equation*}
    x_{ijv} \in \mathbb{Z}_{\geq 0}, \quad \forall i \in I, j \in J, v \in V \quad \text{(restrição de não-negatividade)}
\end{equation*}


\section{Modelo Intermunicipal Com Intermediação Estadual}

Este modelo busca representar o trajeto físico real considerando o intermédio de órgãos estaduais (como a secretaria estadual) no transporte de vacinas entre municípios.

Usaremos
\[ k \in K \]
para representar o centro estadual $k$ dentre os centros estaduais $K$ mapeados.

Com isso, o fluxo pode ser representado da seguinte forma:

\[
\text{Município } i \;\longrightarrow\; \text{Estado } k \;\longrightarrow\; \text{Município } j
\]

\subsection{Variável de Decisão e Parâmetros}

Para o transporte, são necessárias duas etapas de fluxo:

\begin{enumerate}
    \item Município $i \longrightarrow$ Estado $k$ \\
    Para essa etapa, temos:
    \[ x_{ikv} \quad \text{e} \quad C_{ik}\]
    onde:
    \begin{itemize}
    \item $x_{ikv}$ representa a quantidade de vacinas do tipo $v$ transportadas do município $i$ ao centro estadual $k$;
    \item $C_{ik}$ representa o \textbf{custo total} (em horas) para se transportar vacinas do município $i$ ao centro estadual $k$.
    \end{itemize}.
    
    \item Estado $k \longrightarrow$ Município $j$ \\
    Para essa etapa, temos:
    \[y_{kjv} \quad \text{e} \quad C_{kj}\]
    onde:
     \begin{itemize}
    \item $y_{kjv}$ representa a quantidade de vacinas do tipo $v$ transportadas do centro estadual $k$ ao município $j$.
    \item $C_{kj}$ representa o \textbf{custo total} (em horas) para se transportar vacinas do centro estadual $k$ ao município $j$.
    \end{itemize}.
\end{enumerate}

Foram definidas duas variáveis de decisão distintas ( $x_{ikv}$ e $y_{kjv}$ ) para a representação de quantidade de vacinas transportadas.
A separação dessas variáveis é necessária, pois os fluxos de entrada e saída de um centro intermediário podem apresentar distribuições diferentes mesmo que esteja sendo considerado a conservação do fluxo entre etapas. Por exemplo, para um determinado tipo de vacina $v$ em um centro $k$, pode-se ter:

\begin{itemize}
\item Fluxo de entrada:
\[
x_{1kt} = 10, \quad x_{2kt} = 5 \quad \Rightarrow \quad \sum_{i \in I} x_{ikv} = 15
\]

\item Fluxo de saída:
\[
y_{k3t} = 8, \quad y_{k4t} = 7 \quad \Rightarrow \quad \sum_{j \in J} y_{kjv} = 15
\]
\end{itemize}

Observa-se que, embora os fluxos individuais de entrada e saída sejam distintos, o total de vacinas que entra no centro é igual ao total que sai. Dessa forma, a separação entre $x_{ikv}$ e $y_{kjv}$ permite refletir a realidade operacional do sistema logístico.

Para a primeira etapa, as opções de modal podem ser rodoviário ou aéreo. Apesar de ser esperado que o modal escolhido para a segunda etapa seja o rodoviário, deixaremos as opções de escolha dos dois modais e espera-se que o computador escolha sempre o rodoviário para essa situação conforme seria em uma situação real.

De forma análoga ao modelo intermunicipal sem intermediação estadual, temos:

\[
\tau_{ik}^{rod}, \quad \tau_{ik}^{areo}, \tau_{kj}^{rod}, \quad \tau_{kj}^{areo} \in \mathbb{R}_{\geq 0}
\]

onde:
\begin{itemize}
    \item $\tau_{ik}^{rod}$ é o tempo de deslocamento do modal rodoviário entre o município de origem $i$ e o centro estadual de destino $k$.
    \item $\tau_{ik}^{areo}$ é o tempo de deslocamento do modal aéreo entre o município de origem $i$ e o centro estadual de destino $k$.
    \item $\tau_{kj}^{rod}$ é o tempo de deslocamento do modal rodoviário entre o centro estadual de origem $k$ e o município de destino $j$.
    \item $\tau_{kj}^{areo}$ é o tempo de deslocamento do modal aéreo entre o centro estadual de origem $k$ e o município de destino $j$.
\end{itemize}

Sejam:
\begin{itemize}
    \item $C_{ik}^{rod}$, o \textbf{Custo Total do transporte rodoviário da primeira etapa};
    \item $C_{ik}^{areo}$, o \textbf{Custo Total do transporte aéreo  da primeira etapa}.
    \item $C_{kj}^{rod}$, o \textbf{Custo Total do transporte rodoviário da segunda etapa};
    \item $C_{kj}^{areo}$, o \textbf{Custo Total do transporte aéreo  da segunda etapa}.
\end{itemize}


Dessa forma, o \textbf{Custo Total de transporte por modal por etapa} é dado por:

\[
C_{ik}^{rod} = \tau_{ik}^{rod} + \delta^{rod}
\]
\[
C_{ik}^{areo} = \alpha \cdot \tau_{ik}^{areo} + \delta^{areo}
\]
\[
C_{kj}^{rod} = \tau_{kj}^{rod} + \delta^{rod}
\]
\[
C_{kj}^{areo} = \alpha \cdot \tau_{kj}^{areo} + \delta^{areo}
\]


Dessa forma, define-se o \textbf{Custo Total de transporte por Etapa} como:

\[
C_{ik} = \min \left( \quad C_{ik}^{rod} \quad , \quad C_{ik}^{areo} \quad \right),  \quad C_{ik} \in \mathbb{R}_{\geq 0}
\]
\[
C_{kj} = \min \left( \quad C_{kj}^{rod} \quad , \quad C_{kj}^{areo} \quad \right),  \quad C_{kj} \in \mathbb{R}_{\geq 0}
\]

Da mesma forma que no modelo anterior, é necessário considerar que a \textbf{capacidade de transporte} varia de acordo com o modal.

Para a \textbf{primeira etapa} (município $i$ $\rightarrow$ estado $k$), sejam:

\[
u_{ik}^{\text{rod}}, \quad u_{ik}^{\text{areo}} \in \mathbb{R}_{\geq 0}
\]

onde:
\begin{itemize}
    \item $u_{ik}^{\text{rod}}$ representa a capacidade máxima de transporte entre o município $i$ e o centro estadual $k$ via modal rodoviário;
    \item $u_{ik}^{\text{areo}}$ representa a capacidade máxima de transporte entre o município $i$ e o centro estadual $k$ via modal aéreo.
\end{itemize}

A capacidade efetiva da primeira etapa é então definida de forma consistente com o custo escolhido:

\begin{equation}
u_{ik} =
\begin{cases}
u_{ik}^{\text{rod}}, & \text{se } C_{ik}^{\text{rod}} \le C_{ik}^{\text{areo}} \\
u_{ik}^{\text{areo}}, & \text{caso contrário}
\end{cases}
\end{equation}

De forma análoga, para a \textbf{segunda etapa} (estado $k$ $\rightarrow$ município $j$), sejam:

\[
u_{kj}^{\text{rod}}, \quad u_{kj}^{\text{areo}} \in \mathbb{R}_{\geq 0}
\]

onde:
\begin{itemize}
    \item $u_{kj}^{\text{rod}}$ representa a capacidade máxima de transporte entre o centro estadual $k$ e o município $j$ via modal rodoviário;
    \item $u_{kj}^{\text{areo}}$ representa a capacidade máxima de transporte entre o centro estadual $k$ e o município $j$ via modal aéreo.
\end{itemize}

A capacidade efetiva da segunda etapa é definida por:

\begin{equation}
u_{kj} =
\begin{cases}
u_{kj}^{\text{rod}}, & \text{se } C_{kj}^{\text{rod}} \le C_{kj}^{\text{areo}} \\
u_{kj}^{\text{areo}}, & \text{caso contrário}
\end{cases}
\end{equation}

Quanto ao tempo de distribuição de vacinas, o tempo total de entrega de uma dose corresponde, conceitualmente, à soma dos tempos das duas etapas do transporte: do município de origem ao centro estadual e do centro estadual ao município de destino.

No entanto, embora exista conservação de fluxo nos centros estaduais, o modelo não rastreia explicitamente o percurso completo de cada unidade de vacina ao longo das duas etapas. Dessa forma, não é possível garantir que os tempos somados correspondam exatamente ao percurso de uma mesma unidade de vacina.

Ainda assim, como aproximação consistente com a estrutura do modelo e com a organização sequencial das etapas logísticas, o tempo total de distribuição é definido como:

\[
T = \max \left\{ \; C_{ik} \;\middle|\; x_{ikv} > 0 \; \right\} + \max \left\{ \; C_{kj} \;\middle|\; y_{kjv} > 0 \; \right\}
\]

onde:
\begin{itemize}
    \item $T$ representa o tempo total de distribuição;
    \item $C_{ik}$ é o tempo de transporte da primeira etapa;
    \item $C_{kj}$ é o tempo de transporte da segunda etapa.
\end{itemize}

Essa formulação separa as etapa e soma os piores tempos de cada fase e fornece uma estimativa do tempo total necessária para completar a distribuição, mantendo a coerência com a modelagem adotada.


\subsection{Função Objetivo}

Dessa forma, a \textbf{função objetivo} fica definida:

\begin{equation}
    \min z = \sum_{v} \sum_{i} \sum_{k} C_{ik} \cdot x_{ikv} \quad + \quad \sum_{v} \sum_{k} \sum_{j} C_{kj} \cdot y_{kjv}
    \label{eq:intermunicipal_estendido}
\end{equation}

ou, com a penalidade:

\begin{equation}
    \min z = \sum_{v} \sum_{i} \sum_{k} C_{ik} \cdot x_{ikv} \cdot \rho_{iv} \quad + \quad \sum_{v} \sum_{k} \sum_{j} C_{kj} \cdot y_{kjv}
    \label{eq:intermunicipal_estendido}
\end{equation}

onde a primeira parcela representa o custo da primeira etapa do fluxo e a segunda parcela representa o custo da segunda etapa.

A penalidade é aplicada na origem, pois o risco de expiração está associado ao tempo restante no município de origem.

\subsection{Restrições}

A \textbf{restrição de oferta} garante que a quantidade total de vacinas enviadas por um município de origem não exceda o estoque disponível daquele município para cada tipo de vacina.

\begin{equation*}
\sum_{k} x_{ikv} \le o_{ik}, \quad \forall i \in I,\; v \in V
\end{equation*}

A \textbf{restrição de demanda} assegura que toda a demanda por vacinas em cada município de destino seja completamente atendida.

\begin{equation*}
\sum_{k} y_{kjv} = d_{jv}, \quad \forall j \in J, \; v \in V
\end{equation*}

A \textbf{restrição de conservação de fluxo} garante que, para cada centro intermediário, a quantidade total de vacinas recebidas seja igual à quantidade total redistribuída, não havendo acúmulo ou perda de vacinas no centro.

\begin{equation*}
\sum_{i} x_{ikv} = \sum_{j} y_{kjv}, \quad \forall k \in K,\; v \in V
\end{equation*}

A \textbf{primeira restrição de capacidade} é relativa a capacidade de transporte da primeira etapa do fluxo e, nesta restrição, a quantidade de vacinas transportadas não pode ultrapassar esse limite. Logo, temos:
\begin{equation*}
\sum_{v} x_{ikv} \le u_{ik}, \quad \forall i \in I, \; k \in K
\end{equation*}
onde:
\begin{itemize}
    \item $u_{ik}$ representa a capacidade de transporte do veículo na primeira etapa do fluxo.
\end{itemize}

A \textbf{segunda restrição de capacidade} é relativa a capacidade de transporte da segunda etapa do fluxo e, nesta restrição, a quantidade de vacinas transportadas não pode ultrapassar esse limite. Logo, temos:
\begin{equation*}
\sum_{v} y_{kjv} \le u_{kj}, \quad \forall k \in K, \; j \in J
\end{equation*}
onde:
\begin{itemize}
    \item $u_{kj}$ representa a capacidade de transporte do veículo na segunda etapa do fluxo.
\end{itemize}

A \textbf{restrição de viabilidade temporal} é relativa à ideia de remover pares onde o tempo de transporte é maior que o tempo restante até o vencimento. Assim, temos:

\begin{equation*}
    x_{ikv} = 0, \quad \forall i \in I, \; k \in K, \; v \in V \text{ tais que } C_{ik} > e_{iv}
\end{equation*}

A restrição de viabilidade temporal é aplicada apenas na primeira etapa do fluxo. Caso o tempo de transporte entre o município de origem e o centro estadual exceda o tempo restante até a expiração da vacina, o envio é proibido.

Não é necessário impor restrições adicionais sobre os fluxos de saída dos centros estaduais, pois a restrição de conservação de fluxo garante que, na ausência de entrada viável, não haverá saída correspondente, assegurando a consistência do modelo.


A \textbf{restrição de não-negatividade} assegura que as quantidades de vacinas transportadas não podem ser negativas e devem ser números inteiros.
\begin{equation*}
x_{ikv}, y_{kjv} \in \mathbb{Z}_{\geq 0}, \quad \forall i \in I, k \in K, j \in J, v \in V
\end{equation*}

Em suma, o PPLI ficou definido assim:

\begin{equation*}
    \min z = \sum_{v} \sum_{i} \sum_{k} C_{ik} \cdot x_{ikv} \quad + \quad \sum_{v} \sum_{k} \sum_{j} C_{kj} \cdot y_{kjv}
\end{equation*}
sujeito a:
\begin{equation*}
\sum_{k} x_{ikv} \le o_{ik}, \quad \forall i \in I,\; v \in V
\end{equation*}

\begin{equation*}
\sum_{k} y_{kjv} = d_{jv}, \quad \forall j \in J, \; v \in V
\end{equation*}

\begin{equation*}
\sum_{i} x_{ikv} = \sum_{j} y_{kjv}, \quad \forall k \in K,\; v \in V
\end{equation*}

\begin{equation*}
\sum_{v} x_{ikv} \le u_{ik}, \quad \forall i \in I, \; k \in K
\end{equation*}

\begin{equation*}
\sum_{v} y_{kjv} \le u_{kj}, \quad \forall k \in K, \; j \in J
\end{equation*}

\begin{equation*}
    x_{ikv} = 0, \quad \forall i \in I, \; k \in K, \; v \in V \text{ tais que } C_{ik} > e_{iv}
\end{equation*}

\begin{equation*}
x_{ikv}, y_{kjv} \in \mathbb{Z}_{\geq 0}, \quad \forall i \in I, k \in K, j \in J, v \in V
\end{equation*}

\bigskip
